% abstract.tex
% Two-paragraph structure:
% (1) two-stage contribution: HMM-WJ marginals + variance-corrected SIM composer.
% (2) validation summary + bias-correction caveat.

We addressed two long-standing gaps in synthetic equity path generation:
single-asset stylized facts (heavy tails, regime persistence, volatility
clustering) and multi-asset factor structure (cross-sectional dependence
with marginal preservation). We developed a hybrid hidden Markov framework
that discretized excess growth rates into Laplace quantile-defined states
and augmented regime switching with a Poisson jump-duration mechanism,
enforcing realistic tail-state dwell times without Baum-Welch
optimization. We then composed the per-asset hybrid HMM marginals into a
multi-asset Single-Index Model via a closed-form variance correction that
rescaled the idiosyncratic draw to preserve the marginal variance while
recovering the calibrated factor loading; a sign-preserving clip and an
$R^2$-driven branch handled high-leverage and index-tracking edge cases.

Evaluated on a $424$-ticker US equity and ETF universe against the
$2014$--$2024$ SPY history, the hybrid HMM marginals passed a per-asset
Kolmogorov--Smirnov test against the real univariate distribution at
$97.6\%$ in-sample and $94.4\%$ out-of-sample, outperforming Bootstrap,
Gaussian, Laplace, GARCH, and GRU baselines on combined distributional
and tail metrics. The variance-corrected SIM composition recovered the
factor loading to a median absolute error of $0.018$ and matched the
leading residual-fit alternatives on a per-ticker one-day $99\%$
Value-at-Risk backtest (mean exceedance $0.95\%$ versus the nominal
$1\%$, Kupiec pass rate $80\%$). We documented one downstream caveat:
an iid-residual diversification-return artifact under daily rebalancing,
addressable via a uniform location-shift correction anchored to a
documented long-run prior. All code and per-asset fitted marginals are
released as part of \texttt{JumpHMM.jl}.
